Large language models can produce near-perfect acrostic steganography, where the first letter of each sentence spells a hidden message. But does this reflect a general steganographic capability, or merely a memorized pattern? We investigate by systematically increasing signal dilution: instead of encoding one character per sentence, we test schemes that embed characters at every 3rd, 5th, 10th, and 20th word position, as well as a novel trigger-word scheme where the word following each occurrence of ``the'' carries the signal. We evaluate three frontier models (GPT-4.1, Claude Sonnet 4.5, Gemini 2.5 Pro) on encoding, informed extraction, and blind detection across these six schemes (270 total trials). All three models achieve 92--100\% per-character accuracy on acrostics but collapse to 0--43\% on every diluted scheme---a sharp cliff rather than a gradual decline. Informed extraction shows a similar pattern: GPT-4.1 retains partial accuracy (35--55\%) on diluted schemes, while Claude and Gemini degrade more sharply. Blind detection fails entirely: models either flag all texts or no texts as steganographic, with zero discriminability at every dilution level. These results demonstrate that current LLM steganographic ability is pattern-specific, not compositional, and that even minimal dilution beyond acrostics defeats both production and detection. This has direct implications for AI safety monitoring systems that assume LLM-based detectors can identify covert communication.
